% =====================================================================
% Capítulo 4. Marco metodológico
% Estado: 4.1 redactada (06/10/2026); 4.2 a 4.6 pendientes.
% =====================================================================

\chapter{Marco metodológico}\label{cap:metodologia}

Este capítulo describe cómo se ha llevado a cabo el trabajo y por qué se ha
optado por cada alternativa. El capítulo~\ref{cap:marco-teorico} expone en
qué consisten los métodos empleados; aquí se recogen las decisiones concretas
que se han tomado al aplicarlos. El capítulo se organiza en seis partes: las
fuentes de datos y la construcción del conjunto de datos; el protocolo de
comparación equitativa; la especificación de cada modelo; el diseño de la
capa probabilística; las métricas y contrastes empleados; y el entorno de
ejecución y las medidas de reproducibilidad.

\section{Fuentes de datos y construcción del conjunto}\label{sec:datos}

El conjunto de datos del trabajo es una tabla con una fila por hora y
70\,127 filas, que reúne la demanda eléctrica, la temperatura y el calendario
laboral. Todas las transformaciones que conducen a ella están programadas en
los scripts del repositorio, de modo que el conjunto puede regenerarse por
completo a partir de las fuentes públicas originales, como exige el
objetivo~OE1. Esta sección describe cada fuente, las decisiones tomadas al
procesarla y las comprobaciones realizadas, y termina con un resumen de las
limitaciones del conjunto.

% ---------------------------------------------------------------------
\subsection{Visión general y periodo de estudio}\label{sec:datos-vision}

El conjunto cubre, con resolución horaria, el periodo comprendido entre el
1 de enero de 2017 y el 31 de diciembre de 2024 en hora peninsular. Integra
tres fuentes públicas y gratuitas, recogidas en la tabla~\ref{tab:fuentes}:
la demanda real del sistema peninsular, publicada por Red Eléctrica de
España; la temperatura horaria de reanálisis en las 47 capitales de
provincia peninsulares; y el calendario laboral estatal y autonómico
publicado en el Boletín Oficial del Estado. A ellas se suma una fuente
auxiliar, la población provincial del Instituto Nacional de Estadística, que
se emplea para ponderar la temperatura y el calendario.

\begin{table}[htbp]
	\centering
	\caption{Fuentes del conjunto de datos}
	\label{tab:fuentes}
	\small
	\begin{tabularx}{\textwidth}{>{\raggedright\arraybackslash}p{2.5cm} >{\raggedright\arraybackslash}X >{\raggedright\arraybackslash}p{2.8cm} >{\raggedright\arraybackslash}X}
		\hline
		\textbf{Variable} & \textbf{Fuente} & \textbf{Resolución original} & \textbf{Uso en el conjunto} \\
		\hline
		Demanda real peninsular & Red Eléctrica de España, e·sios, indicador 1293 \parencite{ree2026demandareal} & 10~min hasta el 23/05/2022; 5~min después & Variable objetivo, como media horaria \\
		Temperatura del aire a 2~m & Reanálisis ERA5-Land \parencite{munozsabater2021era5land}, servido por Open-Meteo \parencite{zippenfenig2023openmeteo} & Horaria, valor instantáneo; rejilla de 0,1° & 47 series provinciales y media ponderada por población \\
		Población provincial & INE, tabla 2852, Padrón a 1 de enero de 2021 \parencite{ine2021tabla2852} & Anual & Pesos fijos de ponderación \\
		Calendario laboral & Resoluciones anuales de fiestas laborales publicadas en el BOE (sección~\ref{sec:datos-calendario}) & Diaria & Fracción de población en festivo, puentes y tipo de día \\
		\hline
	\end{tabularx}
\end{table}

La elección del periodo responde a tres razones. La primera, expuesta en la
sección~\ref{sec:contexto}, es la coherencia con la aplicación: durante todo
el periodo, la casación del mercado diario, que es la decisión a la que sirve
la previsión estudiada, se realizó en resolución horaria. El argumento se
refiere al mercado diario y no al conjunto del sistema, porque la
liquidación de los desvíos ya pasó a ser cuartohoraria el 1 de diciembre de
2024, dentro del periodo. Tampoco se refiere a la calidad de los datos: la
serie de demanda no cambió con el paso a los productos cuartohorarios; lo
que cambió fue el uso que se da a su previsión. La segunda razón es que el
periodo de prueba debe ser un año natural completo, porque el desglose por
tipo de día del objetivo~OE8 necesita el ciclo anual entero, con su invierno,
su verano, el periodo navideño y los puentes de diciembre; 2024 es el último
año natural completo con mercado diario horario. La tercera es que el
periodo deja fuera el apagón del 28 de abril de 2025, que habría obligado a
tratar como anomalía un episodio excepcional dentro del año de prueba.

El periodo de ocho años es suficiente para los fines del trabajo y mantiene
acotado su coste de cálculo. Proporciona seis años de entrenamiento, de 2017
a 2022, y con ellos seis ciclos anuales completos para estimar la
estacionalidad anual, el ciclo con menos repeticiones en la serie. Al mismo
tiempo, cada año adicional alarga la serie de entrenamiento y encarece el
ajuste de los modelos más costosos, en particular el del SARIMAX, que debe
reajustarse periódicamente durante la evaluación. Se comprobó que la fuente
sirve el año 2017 completo y sin huecos; de no haber sido así, el plan de
reserva era el periodo 2018-2024.

El periodo se divide en tres tramos consecutivos: 2017-2022 para el
entrenamiento, 2023 para la validación, es decir, para elegir la
configuración de cada modelo, y 2024 para la prueba. La partición es
temporal y no aleatoria, porque una partición aleatoria permitiría que un
modelo aprendiera de horas posteriores a las que debe prever. El uso
concreto de cada tramo forma parte del protocolo de comparación, que se
describe más adelante en este capítulo. Todas las decisiones que dependen de
los datos y que se presentan en esta sección, como los umbrales de los
grados-día, se toman exclusivamente con los datos de 2017-2022.

Las variables explicativas se limitan a la temperatura y al calendario.
Ambas cumplen dos condiciones: son públicas y gratuitas, y su valor para el
día siguiente se conoce, o puede preverse, en el momento de emitir la
previsión. El calendario se conoce con exactitud; la temperatura, mediante
una previsión meteorológica, que en este trabajo se sustituye por el
reanálisis, con las consecuencias expuestas en la
sección~\ref{sec:problema}. Otras variables relacionadas con la demanda no
cumplen la segunda condición. El precio del mercado diario para el día
$D+1$, por ejemplo, es el resultado de la casación a la que sirve la propia
previsión, de modo que no se conoce a las 12:00 del día $D$. Otras variables
meteorológicas, como la radiación solar o la humedad, podrían aportar
información adicional, pero se excluyen para acotar el alcance del trabajo.

% ---------------------------------------------------------------------
\subsection{Demanda eléctrica}\label{sec:datos-demanda}

La variable objetivo es la demanda real del sistema eléctrico peninsular,
que Red Eléctrica de España publica como indicador 1293 de su sistema de
información e·sios \parencite{ree2026demandareal}. Se expresa en megavatios
(MW) y tiene resolución inferior a la hora, que además cambia dentro del
periodo: un valor cada diez minutos hasta el 23 de mayo de 2022 a las
19:00~UTC y uno cada cinco minutos a partir de entonces.

Conviene precisar qué recoge esta demanda. Las cifras de demanda que publica
Red Eléctrica no incluyen la energía producida por las instalaciones de
autoconsumo, que se contabiliza aparte: en 2025, tenerla en cuenta habría
elevado la demanda nacional por encima de los 269~TWh
\parencite{ree2026informe2025}.
% Verificado para las estadísticas de REE; para el indicador 1293 es una
% inferencia coherente (la medida del sistema no observa el consumo cubierto
% detrás del contador), no una afirmación explícita de la fuente.
La serie representa, por tanto, la demanda que debe cubrir el sistema, y no
el consumo eléctrico total. La expansión del autoconsumo descrita en la
sección~\ref{sec:contexto} forma parte de su evolución, porque reduce la
demanda observada, sobre todo en las horas centrales del día, sin que ninguna
variable del conjunto de datos la represente. Es una posible fuente de
diferencias entre el periodo de entrenamiento y el de prueba, y se retoma al
discutir los resultados.

Los datos se descargaron a través de la interfaz de programación de e·sios,
que requiere un token personal gratuito, con una petición por mes, 96 en
total. Las peticiones se expresan en tiempo universal coordinado (UTC), que
no tiene horas repetidas ni ausentes; la conversión a hora peninsular se
realiza después, al construir el conjunto de datos
(sección~\ref{sec:datos-tiempo}).

La descarga contiene 558\,030 valores, desde el 1 de enero de 2017 a las
00:00~UTC hasta el 31 de diciembre de 2024 a las 23:55~UTC. Se verificó que
no contiene valores nulos, duplicados ni menores o iguales que cero, y que
la distancia entre dos valores consecutivos es siempre de diez o de cinco
minutos, de modo que la serie no presenta huecos. El único día con un número
de valores distinto del esperado es el 23 de mayo de 2022, el del cambio de
resolución. En consecuencia, no ha sido necesario ningún tratamiento de
huecos.

Para obtener la serie horaria, cada hora $H$ recibe la media de los valores
registrados entre las $H$:00 y las $H$:55; la etiqueta corresponde, por
tanto, al inicio del intervalo. Es la misma convención que siguen los
periodos del mercado diario, en el que el primer periodo va de las 00:00 a
las 01:00. Se emplea la media, y no la suma, a causa del cambio de
resolución: la suma tendría 6 sumandos por hora hasta mayo de 2022 y 12 a
partir de entonces, lo que duplicaría artificialmente la serie en esa fecha,
mientras que la media no depende del número de valores. Al tratarse de la
media de una potencia a lo largo de una hora, el valor obtenido en MW
equivale a la energía consumida en esa hora expresada en MWh.

El número de valores de cada hora se conserva en el conjunto de datos como
variable de control: 47\,251 horas se construyen con 6 valores y 22\,877 con
12, recuentos que cuadran con la descarga, puesto que
$47\,251 \times 6 + 22\,877 \times 12 = 558\,030$. Esas 70\,128 horas
incluyen la última, que en hora peninsular corresponde ya al 1 de enero de
2025 y se descarta, lo que deja las 70\,127 filas del conjunto.

Como consecuencia del cambio de resolución, el periodo de entrenamiento
mezcla ambas resoluciones, mientras que los de validación y prueba se
construyen íntegramente con valores cada cinco minutos. La magnitud medida es
la misma y la media horaria no depende del número de valores, aunque la media
de doce valores puede ser ligeramente más suave que la de seis. Este posible
efecto no se ha cuantificado y se considera menor.

% ---------------------------------------------------------------------
\subsection{Tratamiento del tiempo: hora UTC y hora peninsular}\label{sec:datos-tiempo}

El conjunto de datos conserva la serie indexada en UTC y añade la hora
peninsular como columnas: la hora local completa, la fecha local y la hora
del día. El índice en UTC es continuo, sin horas repetidas ni ausentes, que
es lo que necesitan los modelos que tratan la serie como una sucesión
regular de observaciones, como el SARIMAX y N-HiTS. La hora local, por su
parte, es la que gobierna el comportamiento de la demanda: la actividad se
organiza según el reloj oficial, y el calendario se define por fechas
locales. Por eso, todas las variables que describen el ciclo diario se
construyen sobre la hora local.

Los cambios de hora hacen que los días locales no tengan siempre 24 horas.
El último domingo de marzo, el reloj salta de las 02:00 a las 03:00 y el
día tiene 23 horas; el último domingo de octubre, la hora de 02:00 a 03:00
se repite y el día tiene 25. En el periodo de estudio hay 2905 días locales
de 24 horas, 9 de 23 y 8 de 25. Esos días se mantienen tal como son, con 23
o 25 horas, igual que los periodos que casa el mercado diario en esas
fechas: no se interpola la hora que falta ni se promedian las dos horas que
se repiten. La alternativa de forzar días de 24 horas obligaría a inventar
un dato en marzo y a destruir otro en octubre.

De los nueve días de 23 horas, ocho corresponden a los cambios de hora de
marzo y el noveno es el 1 de enero de 2017, cuya primera hora local, de
00:00 a 01:00, corresponde al 31 de diciembre de 2016 a las 23:00~UTC, que
no forma parte de la descarga. Al pertenecer al periodo de entrenamiento,
esta hora ausente no tiene ningún efecto sobre la evaluación. En el extremo
opuesto, la última hora UTC descargada corresponde ya al 1 de enero de 2025
en hora peninsular y se descarta, como se ha indicado en la
sección~\ref{sec:datos-demanda}.

La coexistencia de ambas escalas de tiempo tiene dos consecuencias para la
modelización. La primera es que el día objetivo de cada previsión puede
tener 23, 24 o 25 horas, de modo que el número de valores que hay que prever
no es constante. La segunda es que los retardos de la demanda, que se
cuentan en horas UTC, no siempre coinciden con la misma hora local del día
anterior o de la semana anterior en las semanas de cambio de hora. La forma
en que cada modelo resuelve ambas cuestiones forma parte del protocolo de
comparación.

% ---------------------------------------------------------------------
\subsection{Temperatura}\label{sec:datos-temperatura}

La temperatura procede del reanálisis ERA5-Land, que reconstruye de forma
coherente la evolución de las variables de superficie terrestre con
resolución horaria desde 1950 y con una resolución espacial nativa de unos
9~km \parencite{munozsabater2021era5land}. Se descarga a través de la
interfaz de datos históricos de Open-Meteo
\parencite{zippenfenig2023openmeteo}, que lo sirve en una rejilla de 0,1°
(unos 11~km), de forma gratuita para uso no comercial y sin registro. La
variable empleada es la temperatura del aire a 2~m. Se descartó el servicio
de datos abiertos de la Agencia Estatal de Meteorología porque no ofrece en
abierto un histórico horario descargable para todo el periodo de estudio.

El modelo de reanálisis se elige de forma explícita. La opción por defecto de
Open-Meteo combina varios modelos sin indicar cuál se usa en cada momento, y
algunos de ellos se generan con la versión del modelo meteorológico vigente
en cada fecha. Un cambio de versión a mitad del periodo podría desplazar la
temperatura y confundirse con un cambio en la demanda. Un reanálisis, en
cambio, se produce con una única versión del modelo para todo el periodo, lo
que garantiza la homogeneidad de la serie. Entre los dos reanálisis
disponibles, ERA5 y ERA5-Land, se elige el segundo por su mayor resolución
espacial.

Se descargan las series horarias de 2017 a 2024 en las capitales de las 47
provincias peninsulares, en UTC, igual que la demanda. Las coordenadas de
cada capital se comprobaron con el servicio de geocodificación de
Open-Meteo, y todas quedan a menos de 10~km de la posición que este
devuelve. Las 47 series tienen 70\,128 horas y ningún valor nulo.

Open-Meteo asigna a cada coordenada una celda de la rejilla según un
criterio configurable. Por defecto, busca una celda situada sobre tierra y
con una elevación similar a la del punto solicitado. Es un criterio
razonable para ponderar por población, porque la población vive en tierra,
pero en las capitales costeras la celda asignada puede no reflejar el efecto
moderador del mar. Todas las celdas asignadas quedan a menos de 11~km de su
capital; la más alejada es la de Cádiz.

Para medir la sensibilidad a esta elección, se comparó ERA5-Land con ERA5
en dos capitales costeras, Barcelona y Cádiz, pidiendo en ERA5 la celda más
próxima, que puede incluir mar. El criterio se fijó antes de ver los
resultados: se mantendría ERA5-Land si la correlación entre las medias
diarias de ambas series era al menos de 0,95 y la diferencia media no
superaba 1~°C. La tabla~\ref{tab:costa} muestra que el criterio se cumple
con holgura en ambas capitales.

\begin{table}[htbp]
	\centering
	\caption{Comparación de ERA5-Land con ERA5 en dos capitales costeras, 2017-2024}
	\label{tab:costa}
	\small
	\begin{tabularx}{0.85\textwidth}{l >{\centering\arraybackslash}X >{\centering\arraybackslash}X}
		\hline
		\textbf{Capital} & \textbf{Correlación de las medias diarias} & \textbf{Diferencia media (ERA5-Land $-$ ERA5)} \\
		\hline
		Barcelona & 0,998 & $-0{,}60$~°C \\
		Cádiz & 0,997 & $+0{,}38$~°C \\
		\hline
	\end{tabularx}
\end{table}

La comparación revela un matiz en Cádiz: la amplitud diaria media, es decir,
la diferencia entre la temperatura máxima y la mínima de cada día, es mayor
en la celda terrestre (8,6~°C) que en la de ERA5 (5,8~°C), lo que refleja el
efecto costero real que la celda terrestre no recoge. Su repercusión en la
media peninsular es reducida, porque Cádiz pesa un 2,8\,\% en la
ponderación por población. Conviene subrayar, además, que ERA5 no es la
temperatura verdadera, sino otro reanálisis: la comparación mide cuánto
cambia el resultado al cambiar de fuente, no cuál de las dos acierta.

La temperatura de Open-Meteo es un valor instantáneo a la hora en punto,
mientras que la demanda de la hora $H$ es una media entre $H$ y $H+1$. Para
que ambas se refieran al mismo intervalo, la hora $H$ recibe la media de las
temperaturas en sus dos extremos:
\begin{equation}
	\tilde{T}_H = \frac{T(H) + T(H+1)}{2},
\end{equation}
que equivale a aproximar la temperatura media de la hora por la regla del
trapecio. La alternativa más sencilla, usar solo $T(H)$, introduciría un
desfase sistemático de media hora. La alineación no añade información
posterior a la que ya supone el escenario de previsión meteorológica
perfecta, porque el instante $H+1$ es el final de la propia hora que se
prevé.

La temperatura peninsular se obtiene como media de las 47 temperaturas
provinciales, ponderada por la población de cada provincia:
\begin{equation}
	T^{\mathrm{pond}}_H = \sum_{p=1}^{47} w_p \, \tilde{T}_{p,H},
	\qquad w_p = \frac{P_p}{\sum_{q} P_q},
\end{equation}
donde $P_p$ es la población de la provincia $p$. La ponderación refleja que
la demanda depende de dónde vive la población: el peso de Madrid es del
15,4\,\% y el de Barcelona del 13,0\,\%, mientras que el de las provincias
menos pobladas no llega al 0,5\,\%. Las poblaciones proceden de las cifras
oficiales del Padrón municipal a 1 de enero de 2021, que suman
43\,869\,377 habitantes en las 47 provincias \parencite{ine2021tabla2852}.

Los pesos son fijos para todo el periodo. Se valoró usar pesos anuales, pero
se descartaron porque introducen una fuga de información pequeña pero real:
la cifra oficial de población a 1 de enero de un año se publica meses
después, de modo que las previsiones de 2024 habrían usado un dato que aún no
existía en el momento de emitirlas. El año de referencia, 2021, pertenece al
periodo de entrenamiento y su cifra se publicó mucho antes del periodo de
prueba. Se eligió 2021, y no 2022, porque la tabla del INE empleada termina
en ese año. Como el reparto de la población entre provincias cambia muy
despacio, el efecto de fijar el año es despreciable.

Además de la temperatura ponderada, el conjunto conserva las 47
temperaturas provinciales con la misma alineación, porque los grados-día se
calculan provincia a provincia (sección~\ref{sec:datos-gradosdia}). Por
último, representar cada provincia por la temperatura de su capital es una
simplificación: en provincias extensas o con relieve marcado, la capital
puede no ser representativa del conjunto de su población.

% ---------------------------------------------------------------------
\subsection{Grados-día de calefacción y de refrigeración}\label{sec:datos-gradosdia}

La relación entre la demanda y la temperatura no es lineal: la demanda
aumenta tanto cuando hace frío como cuando hace calor
(sección~\ref{sec:determinantes}). Para representarla se emplean los
grados-día de calefacción y de refrigeración, que miden cuánto se aleja la
temperatura de una zona neutra. Para la provincia $p$ y el día local $d$,
con temperatura media diaria $T_{p,d}$, se definen como
\begin{equation}
	\mathrm{GDC}_{p,d} = \max\left(T_c - T_{p,d},\, 0\right),
	\qquad
	\mathrm{GDR}_{p,d} = \max\left(T_{p,d} - T_r,\, 0\right),
\end{equation}
donde $T_c$ es el umbral por debajo del cual se activa la calefacción y
$T_r$ el umbral por encima del cual se activa la refrigeración. Los valores
peninsulares se obtienen ponderando por población con los mismos pesos de la
sección~\ref{sec:datos-temperatura}:
$\mathrm{GDC}_d = \sum_p w_p \, \mathrm{GDC}_{p,d}$, y de forma análoga para
la refrigeración.

El orden de las operaciones importa. Los grados-día se calculan primero en
cada provincia y se ponderan después, en lugar de calcularse sobre la
temperatura peninsular ponderada. El motivo es la no linealidad. Supóngase
un día en que la mitad de la población está a 8~°C y la otra mitad a 28~°C.
La temperatura ponderada sería de 18~°C, dentro de la zona neutra, y daría
cero grados-día de ambos tipos, aunque media España esté encendiendo la
calefacción y la otra media el aire acondicionado. Calculados por provincia,
con los umbrales finalmente elegidos ($T_c = 15{,}5$~°C y $T_r = 19{,}5$~°C),
el mismo día tendría $0{,}5 \times 7{,}5 = 3{,}75$ grados-día de calefacción
y $0{,}5 \times 8{,}5 = 4{,}25$ de refrigeración. El ejemplo es extremo, pero
la situación se da de forma atenuada en muchos días de primavera y otoño.

Los umbrales $T_c$ y $T_r$ no se toman de la literatura, sino que se estiman
con los datos, y exclusivamente con los de 2017-2022, para no introducir
información de los periodos de validación y prueba. El procedimiento es el
siguiente:
\begin{enumerate}
	\item Se seleccionan los días laborables de 2017-2022, excluido el periodo
	de la pandemia: 1410 días. Limitarse a los laborables evita que el efecto
	del fin de semana y de los festivos se confunda con el de la temperatura.
	\item Para cada pareja de umbrales candidata se calculan los grados-día
	ponderados y se ajusta por mínimos cuadrados una regresión de la demanda
	media diaria sobre ambos grados-día y una variable indicadora por año. Las
	indicadoras absorben el nivel de demanda de cada año, como el descenso de
	2022, que de otro modo se atribuiría a la temperatura.
	\item Los umbrales se buscan en una rejilla de 0,5~°C, con $T_c$ entre 12
	y 20~°C, $T_r$ entre 16 y 24~°C y la condición $T_c \le T_r$, y se elige la
	pareja con menor suma de cuadrados de los residuos.
\end{enumerate}
Antes de ver el resultado se fijó un criterio de validez: el óptimo no debía
caer en el borde de la rejilla, porque en ese caso el verdadero óptimo podría
estar fuera del intervalo explorado. El resultado se presenta en la
sección~\ref{sec:explo-temperatura}.

Los grados-día no se guardan en el conjunto de datos, sino que se calculan a
partir de las temperaturas provinciales cuando se construyen las variables de
cada modelo. Como los umbrales están fijados con datos de entrenamiento, su
cálculo para cualquier día no requiere información de los periodos de
validación o prueba. La forma en que se combinan con la temperatura horaria
en los modelos es una decisión del protocolo de comparación.

% ---------------------------------------------------------------------
\subsection{Calendario laboral}\label{sec:datos-calendario}

No existe una interfaz oficial que proporcione el calendario laboral de
todas las comunidades autónomas. La fuente empleada son las resoluciones
anuales de la Dirección General de Empleo, y desde la resolución de 2018 de
la Dirección General de Trabajo, que publican en el BOE la relación de fiestas laborales de cada año
\parencite{boe2016fiestas2017,boe2017fiestas2018,boe2018fiestas2019,boe2019fiestas2020,boe2020fiestas2021,boe2021fiestas2022,boe2022fiestas2023,boe2023fiestas2024}.
Se incluyen las fiestas nacionales y las autonómicas de las 15 comunidades
peninsulares, y se excluyen las fiestas locales. La información de las ocho
resoluciones se transcribió a una tabla de referencia con 237 filas, una por
fecha festiva, que indica las comunidades en las que es fiesta.

La tabla se verificó de tres formas. Primero, su coherencia interna: sin
valores nulos, con códigos de comunidad y fechas válidos y sin duplicados.
Segundo, el número de fiestas por comunidad y año, que es de 12 salvo en
nueve casos con 13, todos justificados por las notas de las propias
resoluciones. Tercero, el contraste con la librería \texttt{holidays}
\parencite{holidays2026}, excluidos los domingos: ninguna fecha aparece solo
en la tabla, y ocho aparecen solo en la librería, todas ellas el 3 de
diciembre en Navarra. Esa fecha es una fiesta local común a toda la
Comunidad Foral según su normativa, y el BOE no la recoge como fiesta
autonómica, por lo que se mantiene el criterio de la fuente.

A partir de la tabla se construyen, para cada día local, las variables
siguientes:
\begin{itemize}
	\item \textbf{Fracción de población en festivo} (\texttt{fraccion\_festivo}):
	suma de los pesos de población de las comunidades en las que el día es
	festivo, entre 0 y 1. Los pesos de cada comunidad son la suma de los de sus
	provincias. Es la variable de calendario que usan los modelos, porque
	refleja que un festivo autonómico reduce la demanda en proporción a la
	población afectada. Un festivo nacional vale 1; el 7 de diciembre de 2020,
	festivo en varias comunidades, vale 0,535.
	\item \textbf{Festivo} (\texttt{festivo}): indicador de que la fracción
	anterior es al menos 0,5. Se emplea únicamente para clasificar los días en
	el desglose del objetivo~OE8 y no como entrada de los modelos, por lo que
	no se ajusta con datos ni introduce fuga de información. Tiene un efecto
	frontera conocido: el 9 de diciembre de 2019 queda fuera, con 0,498, y el
	7 de diciembre de 2020 dentro, con 0,535.
	\item \textbf{Puente} (\texttt{puente}): día de lunes a viernes, no
	festivo, situado entre dos días no laborables, entendidos como festivos o
	días de fin de semana. Hay 20 puentes en el periodo de estudio.
	\item \textbf{Tipo de día} (\texttt{tipo\_dia}): clasificación con la
	prioridad festivo, puente, fin de semana y laborable. Un festivo en sábado
	cuenta como festivo. A las tres categorías previstas en el objetivo~OE8 se
	añade la de fin de semana, que de otro modo quedaría mezclada con los
	laborables.
\end{itemize}

La tabla~\ref{tab:tipos-dia} recoge el recuento de días de cada tipo. Los
festivos y los puentes son pocos, lo que deberá tenerse en cuenta al
interpretar el desglose de resultados: el año de prueba tiene 10 festivos y un
solo puente, y las métricas calculadas sobre ellos
serán más inestables que las de los laborables.

\begin{table}[htbp]
	\centering
	\caption{Días de cada tipo en el periodo de estudio, 2017-2024}
	\label{tab:tipos-dia}
	\small
	\begin{tabular}{lr}
		\hline
		\textbf{Tipo de día} & \textbf{Días} \\
		\hline
		Laborable & 1994 \\
		Fin de semana & 826 \\
		Festivo & 82 \\
		Puente & 20 \\
		\hline
		Total & 2922 \\
		\hline
	\end{tabular}
\end{table}

El criterio de la fuente tiene una consecuencia que conviene declarar. Las
resoluciones del BOE no incluyen las fiestas nacionales que caen en domingo,
de modo que esos días aparecen en el conjunto como un domingo corriente, con
fracción de festivo nula. En 2017-2024 hay diez fiestas nacionales de fecha
fija que caen en domingo, verificadas con un script: el 1 de enero de 2017,
el 6 de enero de 2019, el 8 de diciembre de 2019, el 1 de noviembre y el 6 de
diciembre de 2020, el 15 de agosto de 2021, el 1 de mayo y el 25 de
diciembre de 2022, el 1 de enero de 2023, en el periodo de validación, y el
8 de diciembre de 2024, en el de prueba. El efecto esperado es pequeño,
porque la demanda de un domingo ya es baja, pero no nulo: el mínimo horario
del periodo de entrenamiento se registra precisamente el 25 de diciembre de
2022 (sección~\ref{sec:explo-anomalias}). Cuando una fiesta nacional cae en
domingo, varias comunidades la trasladan al lunes, y ese lunes sí figura en
el BOE, como el 7 de diciembre de 2020 citado antes. Estos días se vigilan
en el desglose de resultados.

Quedan fuera dos tipos de días que también alteran la demanda. Las fiestas
locales no se incluyen, ni siquiera las de grandes ciudades como Madrid o
Barcelona, cuyo efecto puede notarse en los errores de esos días. Tampoco se
etiquetan en el conjunto de datos los días especiales que no son festivos,
como el 24 y el 31 de diciembre o el resto del periodo navideño; el 24 y el
31 de diciembre de 2018, por ejemplo, figuran como puentes. El análisis
exploratorio muestra que el periodo navideño concentra los mayores déficits
de demanda (sección~\ref{sec:explo-anomalias}), y por ese motivo el
protocolo de comparación añade una variable específica para ese periodo, con
la evidencia obtenida exclusivamente en 2017-2022.

% ---------------------------------------------------------------------
\subsection{Huecos, valores anómalos y periodo de la pandemia}\label{sec:datos-anomalias}

El conjunto de datos no contiene huecos. La demanda no presenta ninguno en
la descarga (sección~\ref{sec:datos-demanda}), las 47 series de temperatura
están completas y el calendario cubre todos los días del periodo. El
conjunto final tiene un índice horario continuo y ningún valor nulo.

Ningún valor de la demanda se ha corregido ni eliminado. El criterio es que
solo se corrige un dato si es erróneo, y la revisión de los valores extremos
del análisis exploratorio (sección~\ref{sec:explo-anomalias}) no encontró
ninguno: los días con demanda muy alejada de la esperada tienen una
explicación real, ya sea de calendario, como el periodo navideño y el lunes
de Pascua, o meteorológica, como la borrasca Filomena de enero de 2021.
Suavizar esos valores mejoraría el ajuste de los modelos en entrenamiento,
pero los modelos se enfrentarían a episodios del mismo tipo en el periodo de
prueba sin haberlos visto antes.

La única excepción es el periodo de la pandemia. Entre marzo y junio de 2020,
las restricciones de movilidad alteraron la demanda de forma excepcional, y
su inclusión en el entrenamiento enseñaría a los modelos un comportamiento
que no se repite. El conjunto marca esas horas con la columna
\texttt{excluir\_covid}, verdadera entre el 1 de marzo y el 30 de junio de
2020 en fecha local, 2927 horas en total, en lugar de eliminarlas. Borrar las
filas habría dejado un hueco de cuatro meses en el índice, y los retardos que
lo cruzaran habrían unido sin aviso observaciones separadas por ese hueco.
Con la marca, cada modelo aplica la exclusión al construir sus muestras de
entrenamiento, según una regla común que forma parte del protocolo de
comparación. Como el periodo de prueba es 2024, nunca se evalúa sobre esas
horas.

Con lo anterior, la sección cubre los tres elementos que exige el criterio
de cumplimiento de OE1 (tabla~\ref{tab:criterios}): el número de huecos
tratados, que es cero; el criterio aplicado en los cambios de hora
(sección~\ref{sec:datos-tiempo}); y los valores anómalos detectados, que se
describen en la sección~\ref{sec:explo-anomalias}.

% ---------------------------------------------------------------------
\subsection{Congelación y verificación del conjunto}\label{sec:datos-verificacion}

Los datos no se publican en el repositorio, sino que se regeneran con los
scripts a partir de las fuentes. Lo que se congela es, por tanto, el código
que los genera, junto con una huella de contenido que permite comprobar que
el conjunto regenerado es idéntico al empleado en el trabajo. La huella es un
número entero calculado a partir del contenido de todas las celdas de la
tabla: cualquier diferencia en cualquier valor la cambia. Todos los scripts
de análisis y de modelización comprueban la huella antes de empezar y se
detienen si no coincide.

La tabla~\ref{tab:control} recoge las cifras de control del conjunto. Son
cifras estructurales, que describen su forma y no su contenido; los
estadísticos descriptivos de la demanda y de la temperatura se presentan en
el análisis exploratorio, calculados solo sobre el periodo de entrenamiento.
Además de estas cifras, la construcción comprueba que el índice horario es
continuo y sin duplicados, que toda la demanda es positiva y que la tabla
guardada en disco se relee idéntica.

\begin{table}[htbp]
	\centering
	\caption{Cifras de control del conjunto de datos}
	\label{tab:control}
	\small
	\begin{tabular}{lr}
		\hline
		\textbf{Cifra} & \textbf{Valor} \\
		\hline
		Filas (horas) & 70\,127 \\
		Columnas & 60 \\
		Valores nulos & 0 \\
		Horas construidas con 6 valores / con 12 valores & 47\,251 / 22\,877 \\
		Días locales de 24 / 23 / 25 horas & 2905 / 9 / 8 \\
		Horas marcadas como periodo de la pandemia & 2927 \\
		Primera hora local & 01/01/2017 01:00 \\
		Última hora local & 31/12/2024 23:00 \\
		Huella de contenido & 5794786326140387865 \\
		\hline
	\end{tabular}
\end{table}

La huella se calcula con la librería pandas, y su valor depende de que los
datos sean idénticos bit a bit. Durante el trabajo hubo que cambiar la
versión de pandas por un requisito de compatibilidad de otra librería
(sección de entorno de este capítulo), y el conjunto se regeneró con la
nueva versión. La comparación con la versión anterior mostró que 59 de las
60 columnas eran idénticas bit a bit y que la temperatura ponderada difería,
como máximo, en $1{,}4 \times 10^{-14}$~°C. La causa probable, no comprobada
en el código de la librería, es un cambio en el orden de las sumas en coma
flotante al calcular la media ponderada. La copia anterior, leída con la
nueva versión, reproduce la huella original, lo que demuestra que la
diferencia está en los datos y no en el cálculo de la huella. Antes de
comparar se había fijado un criterio de aceptación: diferencias inferiores a
$10^{-10}$ en la temperatura ponderada y el resto de columnas idénticas. Al
cumplirse, se aceptó el conjunto regenerado y su huella pasó a ser la
oficial.

El episodio tiene una consecuencia para la reproducibilidad: el conjunto solo
se reproduce bit a bit con las versiones exactas de las librerías, que por
ese motivo quedan fijadas en el repositorio. Con otras versiones, el conjunto
puede diferir en el orden de la precisión de la máquina y la huella no
coincidiría, aunque los resultados serían a efectos prácticos los mismos.

% ---------------------------------------------------------------------
\subsection{Limitaciones del conjunto de datos}\label{sec:datos-limitaciones}

La tabla~\ref{tab:limitaciones} reúne las limitaciones del conjunto de datos
que se han ido señalando a lo largo de la sección, con su efecto esperado y
el lugar del trabajo en que se retoman.

\begin{table}[htbp]
	\centering
	\caption{Limitaciones del conjunto de datos}
	\label{tab:limitaciones}
	\small
	\begin{tabularx}{\textwidth}{>{\raggedright\arraybackslash}p{4.2cm} >{\raggedright\arraybackslash}X >{\raggedright\arraybackslash}p{3cm}}
		\hline
		\textbf{Limitación} & \textbf{Efecto esperado} & \textbf{Dónde se retoma} \\
		\hline
		Temperatura de reanálisis en lugar de previsión meteorológica & Los resultados son una cota superior del rendimiento en operación & Discusión y conclusiones \\
		Autoconsumo no observado & Posible deriva de la demanda entre el entrenamiento y la prueba, sobre todo en las horas centrales & Discusión \\
		Cada provincia representada por su capital, con celda sobre tierra & Error de representatividad, reducido en la media ponderada & --- \\
		Pesos de población fijos (2021) & Despreciable, por la lentitud del cambio en el reparto de la población & --- \\
		Fiestas locales excluidas & Errores en días concretos de grandes ciudades & Desglose por tipo de día \\
		Fiestas nacionales en domingo sin marcar & Posibles errores el 01/01/2023 (validación) y el 08/12/2024 (prueba) & Desglose por tipo de día \\
		Resolución de la demanda distinta en entrenamiento y prueba & Menor; no cuantificado & --- \\
		Otras variables meteorológicas excluidas & Información potencialmente útil no aprovechada & Conclusiones \\
		\hline
	\end{tabularx}
\end{table}
